\section{Consumption economics: devolver a los desarrolladores el efecto sobre los recursos}

La nube y la IA avanzan hacia el consumption-based pricing. La cuestión no es solo si «cobrar por consumo es barato», sino si la medición puede convertirse en la \term{fitness function} del desarrollador: qué recursos consumió una query, un render, una model call o una data transfer, si eso afecta a otras workloads y si el desarrollador puede optimizar a partir de esa información.

\subsection{Honest metering y recursos atribuibles}

Si las dimensiones de facturación son $u_i$ y los precios unitarios son $p_i$, la factura es
\[
B=\sum_i u_i p_i.
\]
La fórmula en sí no es difícil; la dificultad está en si $u_i$ puede ponerse en correspondencia con las operaciones del desarrollador. Si un join costoso solo se manifiesta como «la base de datos en conjunto se vuelve más lenta y la factura en conjunto sube», la retroalimentación es pobre; solo si el sistema puede señalar el efecto de esa operation en lecturas, CPU, memoria y latencia, el desarrollador puede decidir si agrega un índice, usa caché, divide la consulta o acepta el costo.

\lecturefigure{09-consumption-feedback.png}{La medición del consumo solo genera valor de ingeniería cuando entra en el ciclo de decisiones de desarrollo.}{Subtítulos locales 20:00--22:40; redibujo conceptual.}

\begin{knowledgebox}{Lectura de la figura: la señal de precio debe verse junto con la señal de calidad}
Mostrar solo el costo induce a optimizar en exceso; mostrar solo la latencia oculta el desperdicio de recursos. La interfaz ideal muestra al mismo tiempo request volume, latency, error, cache hit, CPU/memory, transfer y spend, y los asocia con la deployment version. Lo que el desarrollador optimiza es el costo por unidad de resultado de negocio, no la minimización de un solo indicador.
\end{knowledgebox}

\subsection{Noisy workload y aislamiento}

En clase se usó la base de datos como analogía, pero lo que realmente se quería expresar es la workload attribution: en un sistema compartido, una operación costosa puede volver más lentos a sus vecinos sin que exista una medición clara. Estas notas no clasifican de forma absoluta a PostgreSQL o DynamoDB como «sistema malo/bueno», sino que comparan si una configuración concreta puede ofrecer workload isolation, horizontal scaling y operation-level evidence.

\lecturefigure{10-workload-isolation.png}{El aislamiento de recursos y la operation-level accounting hacen explicables el rendimiento y el costo.}{Subtítulos locales 20:30--22:35; redibujo conceptual.}

\begin{warningbox}{El cobro por consumo puede trasladar al usuario la ineficiencia de la plataforma}
Cobrar por consumo no es justo por naturaleza. Si el runtime de la plataforma desperdicia CPU, si la estrategia de cache es ineficiente o si las dimensiones de medición no son explicables, el usuario termina pagando por los problemas de la plataforma. El proveedor debe hacer pública la semántica de la medición, ofrecer recomendaciones de optimización y permitir que el usuario distinga entre el crecimiento del negocio y las regresiones del sistema.
\end{warningbox}

\subsection{Resumen de la sección}

El núcleo de la consumption economics es la attribution. Los recursos, los precios y la calidad del servicio deben mapearse de vuelta a operations y deployments concretos para formar una retroalimentación de ingeniería, en lugar de una factura inexplicable que solo aparece a fin de mes.
